\documentclass[11pt]{article}

\usepackage{amsmath, amsthm, amssymb}
\DeclareMathOperator*{\argmax}{arg\,max}
\DeclareMathOperator*{\argmin}{arg\,min}
\usepackage{array}
\usepackage{forest}
\usepackage[margin=1in]{geometry}

\theoremstyle{definition}
\newtheorem{definition}{Definition}[section]

\theoremstyle{plain}
\newtheorem{theorem}{Theorem}[section]

\title{}
\date{}

\begin{document}

\section{VCG Practice}
Imagine there is one mango ($M$) and one watermelon ($W$) that a social planner is trying allocate. The social planner decides to run a VCG mechanism to sell the fruits. Since there are two different fruits, each agent either gets $\emptyset$ (neither fruit), $\{M\}$ (just the mango), $\{W\}$ (just the watermelon), $\{M,W\}$ (both the mango and the watermelon). Table \ref{tab:valuations} shows each of Alice, Bob, and Charlie's valuation functions for these bundles (they do not have preferences over what others receive):

\begin{table}[h]
\centering
\begin{tabular}{|c|c|c|c|c|}
\hline
 & $\emptyset$ & $\{M\}$ & $\{W\}$ & $\{M,W\}$ \\
\hline
Alice   & 0 & 1 & 3 & 10 \\
\hline
Bob     & 0 & 3 & 6 & 9 \\
\hline
Charlie & 0 & 6 & 1 & 7 \\
\hline
\end{tabular}
\caption{List of agents and their valuations for each bundle.}
\label{tab:valuations}
\end{table}

\noindent \textbf{(a)} Suppose only Alice and Bob participate in the auction. What is the outcome of the VCG mechanism? Explicitly compute the allocation and payment for Alice and Bob.
\\
\\
\noindent \textbf{(b)} Suppose Alice, Bob and Charlie all participate in the auction. What is the outcome of the VCG mechanism? Explicitly compute the allocation and payment for Alice, Bob, and Charlie.
\\
\\
\noindent \textbf{(c)} For the two cases above, what is the total welfare and how much does social planner earn? Explicitly compute the total welfare and the revenue for both cases. 



\section{$k$ Indivisible Units to $k$ Highest Bidders}
Suppose $k< N$. Consider the allocation rule $x:\mathbb{B} \to \mathbb{A}$ which $
\forall b \in \mathbb{B}$ and $\forall i \in [N]$ sets:
$$x_i(b)=\begin{cases}1 & i \in \lambda(b) \\ 0 & \text{else}. \end{cases}$$
where $\lambda(b)$ is the set of $k$ agents such that $\forall i \in \lambda(b),
\forall j \in [N] \backslash \lambda(b)$ we have $b_i \geq b_j$ and if $b_i=b_j$ we
have $i < j$. Note $x(\mathbb{B})$ consists of all tuples in $\mathbb{R}^N_{\geq
0}$ with exactly $k$ entries equal to 1 and the remaining $N-k$ entries equal to
$0$. 
\\
\\
Find a payment rule which implements this allocation rule. Can any other payment rule implement this allocation rule? Justify your answers.




\end{document}
